\chapter{Discussions}

\section{Effectiveness of the Proposed Method}

The proposed latent diffusion-based augmentation method proved to be effective in improving artifact classification performance. The main advantage of this method is its ability to generate realistic, high variance and fully synthetic EEG data rather than introducing small perturbations to existing data as in traditional augmentation techniques. Using this method to extend datasets allows better generalisation and reduces bias toward classes with larger number of datapoints, which is especially relevant for highly imbalanced datasets such as TUAR. 

Extending the set of underrepresented classes balances the overall data distribution, resulting in an increase in the number of training iterations that contain those labels and forcing the classification models to improve their ability to recognise these rare artifact patterns. Interestingly, a similar effect can be observed on classes with higher occurrence in the dataset, where the models also show smaller but consistent improvements in both metrics. The reason behind this behavior is that larger number of training datapoints allows the generative latent diffusion model to provide more realistic samples for dominant classes than for less frequent ones, allowing further improvement on overrepresented classes as well. 

In the context of artifact detection, it is highly positive that the proposed method was able to achieve improvements across all classes. This suggests that the method not only handles underrepresented and overrepresented classes well, but it is also able to generalise to statistically different EEG signal patterns. 

Conducting the diffusion process in a compressed latent space has the benefits of stability and efficiency. This dimensionality reduction results in more stable training compared to other generative models and allows the diffusion model to focus on the relevant information in the input EEG signals.

Overall, the results show that the proposed augmentation method is not only capable of improving EEG artifact classification but also helps overcome data scarcity and class imbalance, demonstrating performance gains in both evaluation metrics.

\section{Augmentation Method Comparison}

The proposed method achieved strong results compared to most other measured augmentations as well. In comparison with the nongenerative methods it outperformed all other augmentations except the segment recombination technique. These results were expected, because generative solutions proved to have stronger effect on performance in the past due to their robustness and ability to introduce significant variability into the dataset. 

Compared to the WGAN-based augmentation, the proposed method achieved similar gains in classification performance. While in some cases the addition of synthetic EEG windows generated by the latent diffusion model reached the highest score, for most models the WGAN-based method outperformed the proposed approach. 

The reason for this could be that diffusion-based approaches often require larger training datasets, while WGANs are able to converge using fewer samples and do not require learning a complex denoising process. Another likely factor is that during the data compression performed by the VAE, some temporal information is lost, which can be critical in case of EEG signals. Diffusion models are also harder to finetune, and if the sampling strategy or the noise scheduling mechanism is suboptimal, the quality of the generated EEG data can drop.

Two existing augmentation pipelines were combined with the proposed method and the WGAN-based generative augmentation. The synthetic EEG windows generated by the diffusion model were able to slightly improve upon the results originally achieved by the combined augmentations. Although the WGAN outperformed the latent diffusion again, the gains were relatively small compared to their standalone improvements over the baseline. The magnitude of the improvement can be explained by the nature of augmentation stacking, where combining augmentations may increase variance, benefitting generalisation, but also adds up the distortions making the datapoints less similar to real data. While too much augmentation can lead to underfitting, the introduction of the proposed augmentation proved to be advantageous by improving both metrics. 

The proposed method was not able to achieve the top score in either experiment, but introduced a generative technique capable of improving existing pipelines.

\section{Limitations}

One of the main limitations of the proposed generative augmentation method is that it requires a separate model for each artifact class. These models are not only trained separately, but for optimal performance separate hyperparameter optimisation are required. This makes it computationally expensive compared to other augmentations that achieve similar gains. 

Another limiting factor is the lack of validation step for the generated datapoints. Although the overall quality of the synthetic EEG windows is measured during the training of the diffusion model, when used, a filtering step could potentially improve the quality of the dataset extension.

While this method helped improve classification metrics on underrepresented major classes, the minor classes with even rarer occurrences are so scarce that the model is unable to learn and generate synthetic extensions. 

Information is potentially lost when the VAE constructs the latent representation of the input signal. While the compression benefits the diffusion process the lost temporal information may result in degraded quality and unrealistic EEG windows.

Another limitation is the dataset specificity. The method was only evaluated and optimised using the TUAR dataset, and while it contains the largest amount of annotated artifacts, testing on other datasets and possibly clinical validation are required to verify its performance in other settings.

\section{Future work}

\begin{description}

\item[\textbf{Generalisation and robustness.}] An important next step is the evaluation on other EEG artifact datasets. Cross-institution analysis using other TUAR versions (e.g., v2) and cross-dataset evaluation would provide information about generalisation potential and real world applicability.

\item[\textbf{Conditional diffusion model.}] Transforming the current method into conditional generation using a single diffusion model could help overcome the limitation of per-class training and optimisation. It would also help reduce the computational cost of both training and synthetic data generation.

\item[\textbf{Improved autoencoders architecture.}] Exploring more advanced models, such as hierarchical VAE or VQ-VAE could help reduce the amount of information lost during the compression. Further developing this part of the architecture has the potential to improve generated EEG signal quality and increase the performance of the proposed generative augmentation method.

\item[\textbf{Data quality evaluation.}] Future work could explore evaluation methods for synthetic EEG data. Investigating artifact specific metrics that describe the similarity of synthetic and real datapoints accurately and validating the generated EEG windows with the help of medical experts could help in further optimising the proposed method.

\end{description}